\section{Magnetic reflection symmetry from co-energy}
\label{sec:symmetry}
\begin{assumption}[Conservative symmetric magnetic map]
At fixed temperature and magnet state, the fundamental or consistently
averaged magnetic relation is single-valued and quasi-static. On a
reflection-invariant current domain it admits a $C^2$ scalar co-energy
$\coenergy(\id,\iq)$ with $\flux=\nabla_{\current}\coenergy$.
The effective geometry, winding distribution, material law, and PM excitation
are invariant under reversal of the q-oriented excitation with the d-oriented
excitation unchanged. Rotor-position harmonics are either neglected or
averaged in a way that preserves this reflection. Hysteretic state changes
and irreversible demagnetization are excluded.
\end{assumption}
The symmetry is an assumption about the effective machine, not a consequence
of its PMSM label. A reflected excitation visits a reflected field solution
with the same scalar co-energy. Therefore
\begin{equation}
 \boxed{\coenergy(\id,\iq)=\coenergy(\id,-\iq).}\label{eq:energy-parity}
\end{equation}
For amplitude-invariant dq variables, physical three-phase co-energy has a
factor $3/2$ in its current differential. Here $\coenergy$ denotes that co-energy
divided by $3/2$, so its gradient is exactly the usual dq flux linkage.
This fixed normalization changes no parity or reciprocity statement.
The connection between construction symmetries, energy descriptions, and
reciprocity is developed in \cite{jebai2014}.

Why does a PM on the d axis not remove this symmetry? It fixes a preferred
\emph{d} direction and can break $i_d\mapsto-i_d$ symmetry. It does not select
a preferred sign of q excitation under the stated effective reflection.
The term $\psi_{\rm pm}i_d$ is a simple example: it is unchanged when $i_q$
changes sign. Bias and d saturation can be more complicated with the same
property. This argument concerns the effective excitation symmetry, not a
naive spatial reflection that ignores the transformation of magnetic vectors.

Differentiating \cref{eq:energy-parity} with respect to $i_d$ gives equality
of d-flux values. Differentiating with respect to $i_q$ gives a minus sign
from the reflected argument. Thus
\begin{equation}
 \boxed{\psid(\id,\iq)=\psid(\id,-\iq),\qquad
 \psiq(\id,\iq)=-\psiq(\id,-\iq).}\label{eq:flux-parity}
\end{equation}
The d component is even; the q component is odd and consequently vanishes at
$i_q=0$. The reflected field can change in magnitude through nonlinear
material response without changing this sign structure.
\begin{figure}[tb]
\centering
\begin{tikzpicture}
\begin{axis}[diagnostic axis,title={Even d flux},ylabel={$\psi_d/\psi_*$},
 xmin=-1.1,xmax=1.1,ymin=.99,ymax=1.04]
 \addplot[strong line,reference quantity,domain=-1:1,samples=81]{1+.02*x^2};
 \addplot[only marks,estimated quantity] coordinates {(-.8,1.0128) (.8,1.0128)};
 \draw[alternative pattern,guide quantity] (axis cs:-.8,1.0128)--(axis cs:.8,1.0128);
\end{axis}
\begin{axis}[diagnostic axis,at={(.5\textwidth,0)},anchor=south west,
 title={Odd q flux},ylabel={$\psi_q/\psi_*$},xmin=-1.1,xmax=1.1,ymin=-1.1,ymax=1.1]
 \addplot[strong line,derived quantity,domain=-1:1,samples=81]{x-.08*x^3};
 \addplot[only marks,estimated quantity] coordinates {(-.8,-.75904) (.8,.75904)};
 \draw[alternative pattern,guide quantity] (axis cs:-.8,-.75904)--(axis cs:.8,.75904);
\end{axis}
\end{tikzpicture}
\caption{A nonlinear symmetric example at $i_d=0$. Reflected q currents
give equal d flux and opposite q flux despite curved responses. Curves are
from the co-energy later specified in \cref{eq:synthetic-energy}; estimated quantity marks
identify a reflected pair.}
\label{fig:magnetic-symmetry}
\end{figure}


\paragraph{What the structural result exposes.}
We now know which channels an ideal map cannot occupy, without knowing its
inductances. We have not determined flux magnitudes or ensured that every
real machine meets the reflection assumption. Broken geometry, unaveraged
harmonics, or path-dependent magnetics invalidate the inference that a
forbidden component must be a measurement error.
